% Content-matching surrogate for the capillary fringe -- paste-ready draft for the WTM paper.
% Voice: A. Wickert conventions (spaced en-dashes " -- " for asides, no em-dashes; active voice).
% NOTATION: reconcile symbols with the paper before use -- here z is depth below the land surface,
% varphi the physical fringe weight, S the numerical spline, psi_a the capillary length.

\subsection{A content-matching surrogate for the capillary fringe}

WTM resolves a single water-table depth per grid cell, so the capillary fringe -- the
tension-saturated zone above the water table -- is a sub-grid feature that the model neither
resolves nor reports. We therefore do not represent the fringe profile; we represent its integrated
effect, and we choose that representation to aid the solver while fixing its physical scale from the
soil.

Let $z \ge 0$ denote depth below the land surface. Let the real capillary taper be a dimensionless
weight $\varphi(z) \in [0,1]$ with $\varphi(0)=1$ at the surface and $\varphi \to 0$ at depth. Its
single physically meaningful scalar is its content,
\begin{equation}
A_{\mathrm{phys}} \;=\; \int_0^{\infty} \varphi(z)\,\mathrm{d}z \;=\; \psi_a ,
\label{eq:content}
\end{equation}
which equals the capillary length $\psi_a$ for any profile normalized to unity at the surface -- a
sharp fringe of height $\psi_a$, an exponential $e^{-z/\psi_a}$, and so on. The profile's own shape
cancels in Eq.~\eqref{eq:content}, and only its content survives.

We replace $\varphi$ by a compact-support spline chosen for numerical conditioning,
\begin{equation}
S(z) = s\!\left(z/w\right), \quad 0 \le z \le w, \qquad S(z)=0 \ \ (z>w),
\end{equation}
with $s:[0,1]\to[0,1]$, $s(0)=1$, $s(1)=0$, monotone and $C^{2}$, and shape factor
\begin{equation}
\kappa_S \;=\; \int_0^{1} s(u)\,\mathrm{d}u .
\end{equation}
We fix the spline width $w$ by requiring that it carry the same content as the physical taper,
\begin{equation}
\int_0^{w} S(z)\,\mathrm{d}z \;=\; A_{\mathrm{phys}}
\qquad\Longrightarrow\qquad
\boxed{\,w \;=\; \psi_a / \kappa_S\,}.
\label{eq:width}
\end{equation}
For the symmetric quintic smoothstep $s(u)=1-(6u^{5}-15u^{4}+10u^{3})$, $\kappa_S=\tfrac12$ and
$w = 2\,\psi_a$.

Equation~\eqref{eq:width} preserves exactly the two quantities that carry physical meaning -- the
taper's content (its zeroth moment, and hence the water it redistributes) and its length scale
$\psi_a$ -- while leaving the pointwise shape $s$ free. Because the content is preserved and the
taper only redistributes water among stores, the domain water balance remains exactly the prescribed
$P-\mathrm{ET}$. The shape is legitimately free: it is sub-grid and unreported, so no model output
and no conserved quantity depends on it. We size the fringe from the soil through a first-order
pedotransfer,
\begin{equation}
\psi_a \;=\; C\,\sqrt{n/K_s},
\label{eq:pedo}
\end{equation}
with $n$ the porosity, $K_s$ the saturated hydraulic conductivity, and $C$ a single calibration
constant, following $K_s \propto r^{2}$ for a characteristic pore radius $r$ and capillary rise
$\psi_a \propto 1/r$.

We do not claim that the spline is the fringe. We collapse the real taper to its one observable,
conserved scalar -- its content $\psi_a$ -- and reproduce that scalar with a solver-friendly spline
by matching areas (Eq.~\ref{eq:width}). The representation is physically scaled and numerically
shaped, and we state both plainly.

% If the effective fringe depth is also wanted, add a first-moment (centroid) constraint
% \int_0^w z\,S\,dz = \int_0^\infty z\,\varphi\,dz, which fixes one shape parameter and trades
% shape freedom for fidelity. The zeroth moment above is the honest minimum for a scalar-wtd model.
